\SourcePage{646}{649}
\section{Scattering at high energies}

If the potential energy is not small compared with $\hbar^2/ma^2$, where $a$
is as usual the range of the field, the energy may nevertheless be so high
that
\begin{equation}
 |U|\ll E\sim\frac{\hbar^2}{ma^2}(ka)^2,
 \tag{131.1}
\end{equation}
while at the same time
\begin{equation}
 |U|\gtrsim\frac{\hbar^2}{ma^2}ka=\frac{\hbar v}{a};
 \tag{131.2}
\end{equation}
it is understood, of course, that
\begin{equation}
 ka\gg1.
 \tag{131.3}
\end{equation}
This is scattering of fast particles for which the Born approximation does not
apply, since neither (126.1) nor (126.2) holds.

\SourcePage{647}{650}
This case can be studied with the wave function (45.9),
\begin{equation}
 \psi=e^{ikz}F(\mathbf r),\qquad
 F(\mathbf r)=\exp\left(-\frac{i}{\hbar v}
 \int_{-\infty}^{z}U\,dz\right),
 \tag{131.4}
\end{equation}
whose applicability requires only $|U|\ll E$.

Section 45 noted that this expression is valid only for $z\ll ka^2$ and hence
cannot be continued directly to distances where (123.3) is already valid. No
such continuation is needed: to calculate the amplitude it is enough to know
the wave function at $a\ll z\ll ka^2$. The integral in the exponent of
$F(\mathbf r)$ may then be extended to infinity:
\begin{equation}
 \psi=e^{ikz}S(\boldsymbol\rho),
 \tag{131.5}
\end{equation}
where
\begin{equation}
 S(\boldsymbol\rho)=\exp[2i\delta(\boldsymbol\rho)],\qquad
 \delta(\boldsymbol\rho)=-\frac1{2\hbar v}
 \int_{-\infty}^{\infty}U\,dz
 \tag{131.6}
\end{equation}
and $\boldsymbol\rho$ is the radius vector in the $xy$ plane.

Fast particles scatter mainly through small angles, which we shall consider.
The momentum change $\hbar\mathbf q$ is then relatively small ($q\ll k$), so
$\mathbf q$ may be regarded as perpendicular to the incident wave vector
$\mathbf k$, lying in the $xy$ plane. The scattered wave is obtained by
subtracting the incident wave $e^{ikz}$, the value of (131.4) at $z=-\infty$,
from (131.5). The amplitude for scattering to
$\mathbf k'=\mathbf k+\mathbf q$ is proportional to the corresponding Fourier
component of the scattered wave\footnote{This method of defining the
scattering amplitude is analogous to that used for Fraunhofer diffraction
(see Vol.~II, \S~61). It is precisely diffraction effects that invalidate
(131.4) at $z\gtrsim ka^2$.}:
\[
 f\sim\int[S(\boldsymbol\rho)-1]e^{-i\mathbf{q\rho}}d^2\rho,
\]
where $d^2\rho=dx\,dy$. The proportionality factor can be found by comparison
with the limiting case of the Born approximation, as below.

\SourcePage{648}{651}
There is another calculation which gives a definite expression directly. Use
(129.2), inserting $\psi$ from (131.4). Since (45.8) gives
\[
 \frac{2m}{\hbar^2}UF=2ik\frac{\partial F}{\partial z},
\]
the amplitude, the coefficient of $e^{ikR_0}/R_0$, is
\[
 f=\frac{ik}{2\pi}\int\frac{\partial F}{\partial z}
 e^{-i\mathbf{q\rho}}dx\,dy\,dz
 =\frac{ik}{2\pi}\int[F(\infty)-F(-\infty)]
 e^{-i\mathbf{q\rho}}dx\,dy.
\]
Substitution of $F$ finally gives\footnote{In two dimensions, for a field
$U(x,z)$, the scattering amplitude defined as in the problem to \S~123 is
\begin{equation}
 f=\sqrt{\frac{ik}{2\pi}}\int[S(x)-1]e^{-iqx}dx.
 \tag{131.7a}
\end{equation}
Here $|f|^2d\varphi$ is the scattering cross-section per unit length along the
$y$ axis, and $\varphi$ is the scattering angle in the $xz$ plane; compare the
Born amplitude in Problem 6 of \S~126.}
\begin{equation}
 f=\frac{ik}{2\pi}\int[S(\boldsymbol\rho)-1]
 e^{-i\mathbf{q\rho}}d^2\rho.
 \tag{131.7}
\end{equation}

If the energy is so high that $\delta\sim|U|a/\hbar v\ll1$, the Born
approximation applies. Expanding $S-1\approx2i\delta$, (131.7) gives, in
agreement with (126.4),
\[
 f=-\frac{m}{2\pi\hbar^2}\int Ue^{-i\mathbf{q\rho}}d^2\rho\,dz.
\]

The total cross-section follows from (131.7) by the optical theorem (125.9).
The forward amplitude is the value at $q=0$, whence
\begin{equation}
 \sigma=2\int\mathop{\rm Re}(1-S)d^2\rho
 =4\int\sin^2\delta(\rho)d^2\rho.
 \tag{131.8}
\end{equation}
The integrand may be regarded as the cross-section for particles with impact
parameters in $d^2\rho$\footnote{Section 152 will generalize (131.7), (131.8)
to scattering by a system of particles.}.

\SourcePage{649}{652}
Formula (131.7) does not assume a central field. It is instructive to obtain it
directly from the exact general formula (123.11) for a central field.

Under (131.1)--(131.3), partial amplitudes with large $l$ dominate. The wave
functions are therefore semiclassical, and (124.1) may be used for $\delta_l$.
Putting $r_0\approx l/k$ and $r^2=z^2+l^2/k^2$ gives
\[
 \delta_l\approx-\frac{m}{\hbar^2}\int_{l/k}^{\infty}
 \frac{U(r)\,dr}{\sqrt{k^2-l^2/r^2}}
 =-\frac{m}{\hbar^2k}\int_0^\infty
 U\left(\sqrt{z^2+l^2/k^2}\right)dz,
\]
which is $\delta(\rho)$ from (131.6) at $\rho=l/k$\footnote{The semiclassical
function $2\hbar\delta(\rho)$ is the change of action produced by $U$ as the
particle passes along its classical trajectory. For a fast particle the
trajectory may be taken as straight, and $2\delta(\rho)$ is then the difference
of the classical action integrals
\[
 \frac1\hbar\int_{-\infty}^{\infty}\sqrt{k^2-2mU/\hbar^2}\,dz
 -\frac1\hbar\int_{-\infty}^{\infty}k\,dz
 \approx-\frac{m}{\hbar^2k}\int_{-\infty}^{\infty}U\,dz.
\]
Thus $2\delta(\rho)$ plays a role analogous to the eikonal of geometrical
optics, and this scattering approximation is often called the \emph{eikonal
approximation}. The amplitude itself is by no means reduced to its
semiclassical expression, since in general the conditions $\theta l\gg1$ and
$\delta_l\gg1$ do not hold.}.

For small angles ($\theta\ll1$), the Legendre polynomials at large $l$ have the
form (49.6)
\[
 P_l(\cos\theta)\approx J_0(\theta l)
 =\frac1{2\pi}\int_0^{2\pi}e^{-i\theta l\cos\varphi}d\varphi.
\]
Insert this in (123.11) and replace the sum over large $l$ by an integral:
\begin{equation}
 f=\frac1k\int\frac{d\varphi}{2\pi}
 \int f_l e^{-i\theta l\cos\varphi}l\,dl
 =\frac{k}{2\pi}\int f_l e^{-i\mathbf{q\rho}}d^2\rho,
 \tag{131.9}
\end{equation}
where $\mathbf q$ and $\boldsymbol\rho$ are two-dimensional vectors of
magnitudes $q=k\theta$ and $\rho=l/k$. Inserting (123.15) with
$\delta_l=\delta(l/k)$ returns (131.7).

\SourcePage{650}{653}
For a central field, after integration over the polar angle $\varphi$ in the
$xy$ plane, with $d^2\rho=\rho\,d\rho\,d\varphi$, (131.7) becomes
\begin{equation}
 f=ik\int_0^\infty\{\exp[2i\delta(\rho)]-1\}
 J_0(q\rho)\rho\,d\rho.
 \tag{131.10}
\end{equation}

Section 126 noted that the Born approximation fails for high-energy scattering
through large angles when the cross-section is exponentially small. The method
developed here also fails under those conditions. Such cases are in fact
semiclassical situations to which perturbation theory does not apply.

By the general rules of the semiclassical approximation (compare \S~52, 53),
the exponent governing the exponential decrease of the cross-section can be
found by considering ``complex trajectories'' in the classically inaccessible
region\footnote{For the pre-exponential factor, see A.~Z.~Patashinskii,
V.~L.~Pokrovskii, and I.~M.~Khalatnikov, ZhETF 45, 989 (1963).}.

In the classical scattering problem, the relation between the deflection angle
$\theta$ in $U(r)$ and the impact parameter $\rho$ is
\begin{equation}
 \frac\pi2\pm\frac\theta2=
 \int_{r_0}^{\infty}\frac{\rho\,dr}
 {r^2\sqrt{1-\rho^2/r^2-U/E}},
 \tag{131.11}
\end{equation}
where $r_0$ is the minimum distance from the center and is a root of
\begin{equation}
 1-\frac{\rho^2}{r^2}-\frac UE=0
 \tag{131.12}
\end{equation}
(see (127.5)). We are interested in angles through which a classical particle
could not be deflected\footnote{This method is not restricted to large $E$;
it applies to every case of exponentially small scattering.}. They correspond
to complex solutions $\rho(\theta)$ of (131.11), with complex $r_0$. From this
function and the classical orbital angular momentum $mv\rho$, calculate
\begin{equation}
 S(\theta)=mv\int\rho(\theta)d\theta.
 \tag{131.13}
\end{equation}

\SourcePage{651}{654}
Here $v$ is the particle's velocity at infinity. The scattering amplitude is
\begin{equation}
 f\sim\exp\left(\frac{iS}{\hbar}\right).
 \tag{131.14}
\end{equation}
Equation (131.12) generally has more than one complex root. The $r_0$ used in
(131.11) must be the one producing the smallest positive value of
$\mathop{\rm Im}S$. In addition, if $U(r)$ has complex singular points, these
too must be included among the competing values of $r_0$\footnote{Recall from
\S~126 that if $U(r)$ has a singularity at real $r$, the cross-section does
not in general decrease exponentially.}.

The region $r\sim r_0$ contributes most to (131.11). At high energies the term
$U/E$ under the radical may then be omitted, and integration gives
\begin{equation}
 \rho=r_0\cos\frac\theta2.
 \tag{131.15}
\end{equation}

If $r_0$ is a singular point of $U(r)$, it depends only on the field, not on
$\rho$ or $E$. Calculating $S$ from (131.13), the amplitude is then
\begin{equation}
 f\sim\exp\left(-\frac{2mv}{\hbar}\mathop{\rm Im}r_0
 \sin\frac\theta2\right).
 \tag{131.16}
\end{equation}

If instead $r_0$ must be a root of (131.12), the exponent depends on the
particular field. For
\[
 U=U_0e^{-(r/a)^2},
\]
which has no singular points at finite distances, the equation
\[
 -\frac UE=1-\frac{\rho^2}{r^2}\approx\sin^2\frac\theta2
\]
gives
\begin{equation}
 r_0=a\sqrt{\ln\left(-\frac{E}{U_0}\sin^2\frac\theta2\right)}.
 \tag{131.17}
\end{equation}
Because $r_0$ depends only weakly on $\theta$, it may be treated as constant in
(131.13); the amplitude is then (131.16) with $r_0$ from (131.17).

\SourcePage{652}{655}
\subsection*{Problems}

\ProblemHeading{1.} Find the total cross-section for scattering by a spherical
square well of radius $a$ and depth $U_0$, under (131.1):
$U_0\ll\hbar^2k^2/m$.

\SolutionHeading We have
\[
 \int_{-\infty}^{\infty}U\,dz=-2U_0\sqrt{a^2-\rho^2}.
\]
By (131.7), the forward amplitude ($q=0$) is
\[
 \begin{aligned}
 f(0)&=-\frac{ik}{2\pi}\int_0^a
 \left[\exp\left(-\frac{2iU_0}{\hbar v}\sqrt{a^2-\rho^2}\right)-1\right]
 2\pi\rho\,d\rho\\
 &=-ika^2\int_0^1(e^{2ivx}-1)x\,dx
 =\frac{ka^2}{2i}\left[1-\frac{i}{v}
 -\frac1{2v^2}(e^{2iv}-1)\right],
 \end{aligned}
\]
where $v=U_0a/\hbar v$ is the ``Born parameter.'' The optical theorem (125.9)
then gives
\[
 \sigma=2\pi a^2\left[1+\frac1{2v^2}-\frac{\sin2v}{v}
 -\frac{\cos2v}{2v^2}\right].
\]

In the Born limit $v\ll1$, this gives $\sigma=2\pi a^2v^2$, in agreement with
Problem 1 of \S~126. In the opposite limit $v\gg1$,
$\sigma=2\pi a^2$, twice the geometrical cross-section. This has a simple
meaning. Every particle with $\rho<a$ is scattered and removed from the
incident beam, so the well behaves as an ``absorbing'' sphere. By Babinet's
principle (see Vol.~II, end of \S~61), the total cross-section is twice the
``absorption'' cross-section.

\ProblemHeading{2.} The same for $U=U_0\exp(-r^2/a^2)$.

\SolutionHeading Here
\[
 \int_{-\infty}^{\infty}U\,dz=a\sqrt\pi U_0e^{-\rho^2/a^2}.
\]
Insertion in (131.7), followed by the evident change of variable, gives
\[
 f(0)=-\frac{ika^2}{2}\int_0^{v\sqrt\pi}(e^{-iu}-1)\frac{du}{u},
\]
where again $v=U_0a/\hbar v$. Thus
\[
 \sigma=2\pi a^2\int_0^{v\sqrt\pi}(1-\cos u)\frac{du}{u}.
\]

\SourcePage{653}{656}
For $v\ll1$, the integrand becomes $u/2$ and
$\sigma=\pi a^2v^2/2$, in agreement with Problem 2 of \S~126 for $ka\gg1$.
For $v\gg1$, write it as $(1-e^{-\lambda u}\cos u)/u$ and then let the small
parameter $\lambda$ tend to zero. Integration by parts gives
\[
 \int_0^{v\sqrt\pi}(1-\cos u)\frac{du}{u}
 \approx\ln(v\sqrt\pi)-\int_0^\infty\ln u\sin u\,du
 =\ln(v\sqrt\pi)+C,
\]
where $C$ is Euler's constant. Hence
\[
 \sigma=2\pi a^2\ln(v\sqrt\pi e^C)\quad\text{for }v\gg1.
\]

\ProblemHeading{3.} Find the small-angle cross-section for electron scattering
by a magnetic field confined to a cylindrical region of radius $a$
(Y.~Aharonov, D.~Bohm, 1959).

\SolutionHeading Let the magnetic field point along the $y$ axis, the axis of
the cylindrical region, and take the incident direction as the $z$ axis. The
scattering is independent of $y$, so the problem below is two-dimensional in
the $xz$ plane.

Outside the cylinder $\mathbf H=0$, but the vector potential is nonzero:
\begin{equation}
 \mathbf A=\frac{\Phi}{2\pi}\Grad{\varphi},
 \tag{1}
\end{equation}
where $\varphi$ is the polar angle in the $xz$ plane and $\Phi$ the magnetic
flux. Indeed, over a circle of radius $r>a$ in this plane,
\[
 \int\mathbf H\,dx\,dz=\oint\mathbf A\,d\mathbf l
 =\frac{\Phi}{2\pi}\int_0^{2\pi}d\varphi=\Phi.
\]
The potential (1) changes the phase of the electron plane wave; by (111.9),
\begin{equation}
 \psi=e^{ikz}\exp\left(-\frac{ie\Phi}{2\pi\hbar c}\varphi\right).
 \tag{2}
\end{equation}
This expression fails in a narrow region along the positive $z$ semiaxis,
since particles passing through the field region have their motion disturbed.
This explains the apparent multivaluedness of (2) on circling the origin, when
$\varphi$ increases by $2\pi$. In reality a discontinuity of finite width lies
near the positive $z$ semiaxis, where (2) is inapplicable; on its two sides
$\varphi$ takes values differing by $2\pi$, for example $\mp\pi$.

For scattering through small $\theta$ with small momentum transfer
$q\approx k\theta$, where $qa\ll1$ and $\theta\ll1$, transverse distances
$x\sim1/q\gg a$ are important and the width of the cut may be neglected. In
the region $z\gg|x|$, the $x$ dependence of $\psi$ on either side

\SourcePage{654}{657}
of the $z$ axis may also be neglected, giving\footnote{Formula (3), like
(131.4), fails at excessively large $z$, when diffraction effects become
important.}
\begin{equation}
 \psi=e^{ikz}F(x),\qquad
 F(x)=\begin{cases}
 \exp(-ie\Phi/2\hbar c),&x>0,\\
 \exp(ie\Phi/2\hbar c),&x<0.
 \end{cases}
 \tag{3}
\end{equation}

The ``two-dimensional'' amplitude is calculated from (131.7a)\footnote{For
$q\ne0$, this formula can also be obtained, as explained in the text, without
using the Schrödinger equation in a potential field. We denote the scattering
angle by $\theta$ to distinguish it from the polar angle $\varphi$.}. For
$q\ne0$,
\[
 f=\sqrt{\frac{ik}{2\pi}}\left\{
 \left[\exp\left(-\frac{ie\Phi}{2\hbar c}\right)-1\right]
 \int_0^\infty e^{-iqx}dx+\text{c.c.}\right\}.
\]
Evaluate the integral by inserting $e^{-\lambda x}$ and then taking
$\lambda\to0$. The result is
\[
 f=\frac{i}{q}\sqrt{\frac{2ik}{\pi}}
 \sin\frac{e\Phi}{2\hbar c}.
\]
Hence
\begin{equation}
 d\sigma=|f|^2d\theta=\frac{2}{\pi k\theta^2}
 \sin^2\frac{e\Phi}{2\hbar c}\,d\theta
 \tag{4}
\end{equation}
and for $e\Phi/\hbar c\ll1$,
\[
 d\sigma=\frac{e^2\Phi^2}{2\pi k\hbar^2c^2\theta^2}d\theta,
\]
the perturbative result.

Notice the periodic dependence of (4) on the magnetic-field strength and the
divergence of the total cross-section from $\theta\to0$, even though the field
is confined to a finite region. Both are specifically quantum effects.
